\textbf{1.} Tout triangle peut être placé (par une isométrie) de sorte que
$A$ soit l'origine, $B$ sur la demi-droite des abscisses positives, donc
$B(c\,;0)$, et $C$ dans le demi-plan $y > 0$. Si $C$ a pour coordonnées
$(x_C\,;h)$ avec $h > 0$, il suffit de poser $x = \frac{c}{2} - x_C \in \mathbb{R}$.
Alors
\[
\begin{aligned}
b^2 &= AC^2 = \Bigl(\frac{c}{2} - x\Bigr)^2 + h^2, \\
a^2 &= BC^2 = \Bigl(c - \frac{c}{2} + x\Bigr)^2 + h^2 = \Bigl(\frac{c}{2} + x\Bigr)^2 + h^2,
\end{aligned}
\]
et, la hauteur issue de $C$ valant $h$ (distance de $C$ à l'axe des
abscisses), $\mathcal{S} = \dfrac{ch}{2}$.

\textbf{2.} Le triangle est équilatéral si et seulement si $C$ est le point
$\left(\frac{c}{2}\,;\frac{\sqrt3}{2}c\right)$, c'est-à-dire $x = 0$ et
$h = \frac{\sqrt3}{2}c$, soit $x = y = 0$.

\textbf{3.} On a $\left(\frac{c}{2} - x\right)^2 + \left(\frac{c}{2} + x\right)^2
= \frac{c^2}{2} + 2x^2$ et $4\sqrt3\,\mathcal{S} = 2\sqrt3\,ch$. Donc
\[
\begin{aligned}
a^2 + b^2 + c^2 - 4\sqrt3\,\mathcal{S}
&= \frac{c^2}{2} + 2x^2 + 2h^2 + c^2 - 2\sqrt3\,ch \\
&= 2x^2 + 2\Bigl(h^2 - \sqrt3\,ch + \frac{3c^2}{4}\Bigr) \\
&= 2x^2 + 2\Bigl(h - \frac{\sqrt3}{2}c\Bigr)^2 ,
\end{aligned}
\]
c'est-à-dire $a^2 + b^2 + c^2 - 4\sqrt3\,\mathcal{S} = 2x^2 + 2y^2$.

\textbf{4.} Une somme de carrés est positive, donc
$a^2 + b^2 + c^2 \geqslant 4\sqrt3\,\mathcal{S}$. L'égalité a lieu si et
seulement si $x = y = 0$, c'est-à-dire (question 2) si et seulement si le
triangle est équilatéral.

\textbf{5.} Le triangle $3$, $4$, $5$ est rectangle ($9 + 16 = 25$), d'aire
$\mathcal{S} = \frac{3 \cdot 4}{2} = 6$. On a $a^2 + b^2 + c^2 = 50$ et
$4\sqrt3 \cdot 6 = 24\sqrt3$. Or $50 > 24\sqrt3 \iff 2500 > 1728$, ce qui est
vrai.

\medskip
\textbf{Remarque (seconde preuve, plus géométrique).}
Soit $C'$ le sommet du triangle équilatéral $ABC'$ construit sur $[AB]$ du
côté de $C$, et $p = CC'$. Dans le triangle $CBC'$, on a $BC = a$, $BC' = c$
et $\widehat{CBC'} = |\beta - 60^\circ|$, où $\beta = \widehat{ABC}$. La loi
des cosinus donne
\[
p^2 = a^2 + c^2 - 2ac\cos(\beta - 60^\circ)
= a^2 + c^2 - ac\cos\beta - \sqrt3\,ac\sin\beta .
\]
Or $ac\sin\beta = 2\mathcal{S}$ et $2ac\cos\beta = a^2 + c^2 - b^2$, d'où
\[
p^2 = a^2 + c^2 - \frac{a^2 + c^2 - b^2}{2} - 2\sqrt3\,\mathcal{S}
= \frac{a^2 + b^2 + c^2 - 4\sqrt3\,\mathcal{S}}{2} \geqslant 0,
\]
avec égalité si et seulement si $C = C'$. \hfill $\blacksquare$
